% id: 131-25
% book: 베이직쎈 확률과 통계 · 07 통계적 추정
% section: 유형 082 표본평균의 확률; 표본의 크기 구하기
% figure: tikz:fig-131-25
% answer: 
% answer_source: 
% variant_level: 0 (원본 전사)
% 이 파일은 build.mjs 가 items.json 에서 만든다. 고칠 때는 items.json 을 고친다.
\smash{\raisebox{1.5mm}{\dmkichul{베이직쎈 확률과 통계 131쪽 25번}}}
\noindent\begin{minipage}[t]{0.58\linewidth}\vspace{0pt}\raggedright 정규분포 $\mathrm{N}(200{,}\mkern3mu\mathopen{}60^2)$을 따르는 모집단에서 크기가 $n$인 표본을 임의추출할 때, 표본평균 $\overline{X}$에 대하여 $\mathrm{P}(\overline{X}\le 170)=0.07$이다. 다음에 답하시오.\end{minipage}\hfill\begin{minipage}[t]{0.4\linewidth}\vspace{0pt}\centering \resizebox{\ifdim\width>\linewidth\linewidth\else\width\fi}{!}{% 131-25(인쇄 131쪽) — 소문항 ⑵ 오른쪽 표준정규분포표. 스캔 표라 tabular 로 다시 조판(칸 값 원문 그대로).
{\setlength{\tabcolsep}{4pt}\renewcommand{\arraystretch}{1.3}\small%
\begin{tabular}{|c|c|}
\hline
$z$ & $\mathrm{P}(0\le Z\le z)$ \\
\hline
$1.0$ & $0.34$ \\
\hline
$1.5$ & $0.43$ \\
\hline
$2.0$ & $0.48$ \\
\hline
\end{tabular}}
}\end{minipage}\par
\dmsubprob{1} 확률변수 $Z$가 표준정규분포 $\mathrm{N}(0{,}\mkern3mu\mathopen{}1)$을 따를 때, $\mathrm{P}(\overline{X}\le 170)=\mathrm{P}(Z\le a)$가 되도록 하는 $a$를 $n$에 대한 식으로 나타내시오.
\dmsubprob{2} 오른쪽 표준정규분포표를 이용하여 $n$의 값을 구하시오.
